\chapter{Principes de dénombrement}
\label{workbook:chapter:07}
\section{Le principe fondamental du dénombrement}
\label{workbook:section:07.01}
\begin{exerciseblock}{Le principe fondamental du dénombrement}
\item Un menu propose $4$ entrées, $6$ plats et $3$ desserts. Combien de repas complets sont possibles? Combien sans dessert?
\item Un code comporte deux lettres suivies de trois chiffres, avec répétition permise. Combien de codes existe-t-il?
\item Dessiner un arbre pour deux choix successifs de tailles $3$ et $4$, puis relier le nombre de feuilles au produit.
\end{exerciseblock}
\section{Factorielle et permutations}
\label{workbook:section:07.02}
\begin{exerciseblock}{Factorielle et permutations}
\item Calculer $8!/(6!)$ sans développer toutes les factorielles.
\item De combien de façons peut-on placer $7$ livres distincts sur une étagère? Et si deux livres précis doivent rester ensemble?
\item Combien d'anagrammes distinctes du mot MATHEMATIQUE peut-on former? Identifier les répétitions.
\end{exerciseblock}
\section{Arrangements}
\label{workbook:section:07.03}
\begin{exerciseblock}{Arrangements}
\item Parmi $12$ personnes, choisir un président, un vice-président et un secrétaire. Combien de possibilités?
\item Calculer $A_{10}^4$ et expliquer chaque facteur du produit.
\item Un identifiant utilise $4$ lettres distinctes choisies parmi $8$. Comparer le nombre d'identifiants avec et sans répétition.
\end{exerciseblock}
\section{Combinaisons}
\label{workbook:section:07.04}
\begin{exerciseblock}{Combinaisons}
\item Choisir $5$ personnes parmi $13$ : calculer $\binom{13}{5}$.
\item Expliquer combinatoirement l'identité $\binom nk=\binom n{n-k}$.
\item Une équipe de $6$ personnes doit compter exactement $2$ membres parmi $5$ spécialistes et $4$ membres parmi $9$ généralistes. Combien d'équipes?
\end{exerciseblock}
\section{Critères de sélection d'une méthode}
\label{workbook:section:07.05}
\begin{exerciseblock}{Critères de sélection d'une méthode}
\item Pour chacune des situations, choisir entre produit, permutation, arrangement ou combinaison : former un mot de passe; classer tous les finalistes; choisir un comité; attribuer trois médailles.
\item Construire un arbre de décision utilisant les questions : tous les objets? ordre important? répétition permise?
\item Deux méthodes donnent $\binom{10}{3}3!$ et $A_{10}^3$. Expliquer pourquoi elles coïncident.
\end{exerciseblock}
\section{Comptage par complément et inclusion-exclusion}
\label{workbook:section:07.06}
\begin{exerciseblock}{Comptage par complément et inclusion-exclusion}
\item Combien de chaînes binaires de longueur $8$ contiennent au moins un $1$?
\item Dans une classe de $40$, $23$ étudient l'anglais, $18$ l'espagnol et $9$ les deux. Combien étudient au moins une langue? aucune?
\item Compter les entiers de $1$ à $100$ divisibles par $2$ ou par $5$.
\end{exerciseblock}
\section{Voir les choix comme un arbre}
\label{workbook:section:07.07}
\begin{exerciseblock}{Voir les choix comme un arbre}
\item Construire l'arbre des résultats de trois lancers d'une pièce et compter les chemins donnant exactement deux faces.
\item Un trajet comporte trois embranchements : $2$, puis $3$, puis $2$ choix. Représenter et compter les trajets.
\item Expliquer comment un arbre permet de repérer un oubli ou un double comptage.
\end{exerciseblock}
\section{Pourquoi l'ordre change le comptage}
\label{workbook:section:07.08}
\begin{exerciseblock}{Pourquoi l'ordre change le comptage}
\item Comparer le nombre de podiums de trois personnes parmi $10$ au nombre de comités de trois personnes parmi $10$.
\item Lister les ordres possibles d'un même groupe $\{A,B,C\}$ et expliquer le facteur $3!$.
\item Un tirage de $4$ cartes est-il ordonné dans une main de poker? dans une suite de cartes révélées? Calculer les deux nombres à partir de $52$.
\end{exerciseblock}
\section{Construire un comptage fiable}
\label{workbook:section:07.09}
\begin{exerciseblock}{Construire un comptage fiable}
\item Compter les nombres à quatre chiffres distincts formés avec $0,1,2,3,4,5$, sans zéro initial.
\item Combien de mots de longueur $6$ sur l'alphabet $\{A,B,C\}$ utilisent les trois lettres au moins une fois?
\item Donner deux contrôles de plausibilité pour un résultat de dénombrement : borne supérieure et vérification sur un petit cas.
\end{exerciseblock}
\section*{Synthèse du chapitre}
\begin{synthesisbox}
\item Combien de mots distincts peut-on former avec les lettres de PROBABILITE? Parmi eux, combien commencent par une voyelle?
\item Une association choisit $5$ personnes parmi $8$ femmes et $7$ hommes, avec au moins $2$ femmes et au moins $2$ hommes. Compter.
\item Combien de codes de $6$ chiffres contiennent exactement deux zéros, ne commencent pas par zéro et n'utilisent aucun autre chiffre deux fois?
\end{synthesisbox}
